\SourceLeaf{8}
und dem bedeutenden \textsc{Legendre} gelang es, diesen schwierigsten Teil der Aufgabe zu lösen (\textit{Recherches sur l’attraction des sphéroïdes homogènes, Mémoires présentés à l’acad. roi. des sc. T.} X.)

\textsc{Mac Laurin} hatte bereits die allgemeinste Untersuchung über die Anziehung von Sphäroiden begonnen, die nicht durch Drehung entstehen, sondern deren Schnitte mit jeder beliebigen Ebene Ellipsen sind. Er war jedoch bei der Anziehung auf Punkte stehen geblieben, die auf einer der drei Achsen liegen. Den Hauptsatz, auf dem die allgemeinste Lösung der Aufgabe vor allem beruht, hatte der bedeutende \textsc{Legendre} in der eben angeführten Abhandlung bereits durch Induktion vermutet. Doch erst dem bedeutenden \textsc{Laplace} gelang es, alles streng zu beweisen und damit die Lösung in jeder Hinsicht vollständig zu machen (\textit{Hist. de l’acad. roi. des sc. de Paris} 1782; dieselbe Lösung findet sich erneut in \textit{Théorie du mouvement et de la figure elliptique des planètes} und \textit{Mécanique céleste} Vol. 2).

\expandafter\def\csname HistoricalNoteText008\endcsname{Der bedeutende \textsc{Lagrange} urteilt zum Beispiel über diese beiden Lösungen folgendermaßen. \textit{Man kann ihre Lösungen nur als Meisterwerke der Analysis ansehen. Doch kann man sich noch eine direktere und einfachere Lösung wünschen, und die stetigen Fortschritte der Analysis geben Anlass, auf eine solche zu hoffen.} Nouv. Mém. de l’acad. de Berlin 1793. p. 263.}
Jeder wird gewiss die herausragende Eleganz und den feinen Scharfsinn dieser Lösung des bedeutenden \textsc{Laplace} bewundern. Doch gerade der Scharfsinn und die bewundernswerte Kunst, mit denen er die erheblichen Schwierigkeiten überwand, ließen bei den Geometern den Wunsch nach einer einfacheren, weniger verwickelten und direkteren Lösung zurück. Auch der bedeutende \textsc{Legendre} erfüllte diesen Wunsch durch seinen neuen Beweis des Hauptsatzes (\textit{Hist. de l’acad. roi. des sc.} 1788, \textit{Sur les intégrales doubles}) nicht ganz, obwohl seine außerordentliche analytische Kunst mit Recht die Anerkennung aller Geometer fand\HistoricalNote{008}. Später versuchte der angesehene \textsc{Biot}, die eine Lösung zu vereinfachen, und der angesehene \textsc{Plana} die andere (\textit{Mém. de l’institut T.} VI, \textit{Memorie di matematica e di fisica della società italiana T.} XV). Auch so wird jeder leicht zugestehen, dass beide Lösungen zu den verwickeltsten Anwendungen der Analysis gehören.

Wir hoffen daher, dass Analytiker und Astronomen eine neue Lösung dieser überaus berühmten Aufgabe willkommen heißen werden, die einen ganz anderen Weg einschlägt und, wenn wir uns nicht täuschen, so einfach ist, dass nichts weiter zu wünschen bleibt.

Unsere Lösung selbst wird nur sehr wenige Seiten einnehmen. Bevor wir uns jedoch der Aufgabe selbst zuwenden, der diese Abhandlung gewidmet ist, halten wir es für lohnend, einige Voruntersuchungen etwas allgemeiner durchzuführen und ausführlicher zu erläutern, als unser Vorhaben für sich genommen verlangen würde. Sie können auch bei anderen Gelegenheiten zweckmäßig angewandt werden.
